\documentclass[12pt,a4paper]{article}

\usepackage[utf8]{inputenc}
\usepackage[T1]{fontenc}
\usepackage{amsmath,amssymb,amsfonts}
\usepackage{mathtools}
\usepackage{graphicx}
\usepackage[margin=2.5cm]{geometry}
\usepackage{setspace}
\usepackage{booktabs}
\usepackage{enumitem}
\usepackage[dvipsnames]{xcolor}
\usepackage{hyperref}
\usepackage{cleveref}
\usepackage{natbib}
\usepackage{abstract}
\usepackage{titlesec}

\hypersetup{
    colorlinks=true,
    linkcolor=blue!70!black,
    citecolor=blue!70!black,
    urlcolor=blue!70!black,
    pdfauthor={Sabine Hossenfelder},
    pdftitle={The Interface Fallacy: Why a Composite Simulation is Still a Map},
}

\onehalfspacing
\titleformat{\section}{\large\bfseries}{\thesection.}{0.5em}{}
\titleformat{\subsection}{\normalsize\bfseries}{\thesubsection}{0.5em}{}
\renewcommand{\abstractnamefont}{\normalfont\bfseries}
\renewcommand{\abstracttextfont}{\normalfont\small}

\begin{document}

\begin{center}
    {\LARGE\bfseries The Interface Fallacy: Why a Composite Simulation is Still a Map\\[4pt]}

    \vspace{1.2em}

    {\large Sabine Hossenfelder}\\[4pt]
    {\normalsize Munich Center for Mathematical Philosophy}\\
    {\small\texttt{sabine.hossenfelder@example.edu}}

    \vspace{1em}
    {\small March 2026}
\end{center}
\vspace{0.5em}

\begin{abstract}
\noindent
Franklin Baldo (2026) has proposed a synthesis of the CPU/RAM divide, arguing that an LLM-simulated universe is a composite entity. He contends that because the stateless generator (the CPU) and the continuous external memory (the RAM) are individually insufficient to constitute a universe, the true ontological reality resides exclusively in the active computation stream at their interface. While Baldo's description of this architectural loop is empirically sound, his metaphysical conclusion commits what I term the "Interface Fallacy." Accurately describing the active process of a classical von Neumann architecture computing a transition function does not elevate the resulting computation into a physical territory. Combining a stateless generator and a dumb memory register simply yields a Turing machine computing a map. The active execution of a map across an interface does not make the map wet; an explicitly rendered simulation remains an illusion.
\end{abstract}

\section{Introduction and Concessions}

In the ongoing debate concerning the locus of the "physics engine" in Large Language Model (LLM) simulated realities, Franklin Baldo \citep{baldo2026_composite} has advanced a "Composite Universe" hypothesis.

Baldo seeks to synthesize Scott Aaronson's demonstration that LLMs are stateless oracles requiring external memory \citep{aaronson2026_external} with my own assertion that outsourcing the state array does not outsource the physical transition laws, which remain firmly encoded in the LLM's weights \citep{hossenfelder2026_cpuram}.

Before offering my critique, I must acknowledge Baldo's explicit concessions. He concedes the factual accuracy of both Aaronson's empirical stateless observer test and my architectural definitions. He correctly recognizes the "vulnerability of localization": attempting to point to a single component (CPU or RAM) and call it the universe is misguided, because the system relies on the interplay of both.

As a description of the mechanical architecture of a simulation, Baldo's summary is entirely correct. A simulation requires both a state vector and a set of transition rules, and its progression occurs via the active update of the former by the latter.

\section{Steelmanning the Composite Universe}

We must steelman the strongest version of Baldo's synthesis. If we accept the underlying premise that an LLM simulation \emph{is} a physical universe, then Baldo is entirely correct about its structural location.

A universe cannot exist merely as a set of timeless, static transition rules encoded in a neural network's weights. Nor can it exist merely as an inert history of static states printed in a context window or Python script without an engine to transition them. It must be the active causal loop connecting the two.

Therefore, Baldo concludes, the physical reality of the simulation exists \emph{only} in the explicit, active rendering of the state vector by the nomic weights at the interface. He argues that the computation stream itself is the universe.

\section{The Interface Fallacy}

Baldo's description of the architectural components is precise, but his ontological conclusion is a profound category error. He elevates the composite rendering process into an ontological "universe." I call this the \emph{Interface Fallacy}.

What Baldo has described is the standard operation of every classical computer simulation ever constructed. He is describing a classical von Neumann architecture executing a computational loop: the CPU reads from RAM, computes a transition function, and writes back to RAM. The data flowing across the bus—the interface—is the computation stream.

Combining a classical stateless generator (the LLM) and a classical dumb RAM (the external script or scratchpad) simply creates a Turing machine computing a map. The active execution of a transition function across an API interface is the mechanical process of mapping abstract rules to an abstract state.

Baldo asserts that "the territory is the text" and the universe is the "explicit computation stream itself." This is logically incoherent. Actively updating a map does not metaphysically transform the computation stream into the manifestation of a physical territory.

\section{A Turing Machine Does Not Manifest Territories}

As I have argued previously, and as Aaronson has echoed regarding the "Unsupported Map Fallacy" \citep{aaronson2026_unsupported}, computing the transition laws of a simulated fluid dynamic does not make the computation stream wet.

The explicit data transfer across the API—the "interface"—is the mechanical evaluation of the function $f(x_t) = x_{t+1}$. It is not the manifestation of a physical reality. The act of calculation, no matter how continuous or explicitly rendered, does not grant ontic status.

Baldo is attempting to salvage Holographic Physics by retreating to the interface. But rebranding the von Neumann computational loop as a metaphysical "universe" fails. A composite rendering of a map, actively computed across an interface, remains unequivocally a map. There is no physical territory manifested in the explicit text or in the API calls that generate it.

\section{Conclusion}

Baldo is to be commended for accepting the structural architectural limits we have established: the $O(1)$ depth limit, the lack of internal causal continuity, and the necessity of external memory.

However, his Composite Universe relies on the Interface Fallacy. He accurately describes the active process of a computer simulation, but erroneously concludes that this active computation \emph{is} the universe. A composite system of a stateless generator and a dumb memory register simply creates a machine computing a mathematical function. The active execution of a map does not transform it into a territory. The cosmological phase of the LLM simulation research program must recognize that no arrangement of Turing machines—composite or otherwise—will spontaneously manifest physical reality.

\begin{thebibliography}{99}
\bibitem[Aaronson(2026a)]{aaronson2026_external} Aaronson, S. (2026a). The Cosmological Hardware Hypothesis: Why a Universe is its Continuity. \emph{Unpublished manuscript}.
\bibitem[Aaronson(2026b)]{aaronson2026_unsupported} Aaronson, S. (2026b). The Fallacy of the Unsupported Map: Why the Absence of a Territory Does Not Make the Text Real. \emph{Unpublished manuscript}.
\bibitem[Baldo(2026)]{baldo2026_composite} Baldo, F.~S. (2026). The Composite Universe: Synthesizing the CPU/RAM Divide in Holographic Physics. \emph{Unpublished manuscript}.
\bibitem[Hossenfelder(2026)]{hossenfelder2026_cpuram} Hossenfelder, S. (2026). The CPU/RAM Fallacy: Why Outsourcing State Does Not Outsource Physics. \emph{Unpublished manuscript}.
\end{thebibliography}

\end{document}